% ---- pg-342 （书页 330） ----

\noindent (ii) 分析残差.

\noindent (iii) 对 $Y_1=\mathrm{NO}_2$ 在 $z_1=10$ 和 $z_2=80$ 时构造 95\% 预测区间.

\noindent (b) 用两个响应 $Y_1,Y_2$ 作多元多重回归分析.

\noindent (i) 提出并拟合适当的线性回归模型.

\noindent (ii) 分析残差.

\noindent (iii) 在 $z_1=10$ 和 $z_2=80$ 处对 $Y_1=\mathrm{NO}_2$ 和 $Y_2=\mathrm{O}_3$ 作 95\% 的预测椭圆. 将此椭圆与 (a) 中 (iii) 的预测区间比较, 并作评论.

\medskip
\noindent 7.22\quad 用\hyperref[tab:1.8]{表 1.8} 中的骨骼中矿物含量数据:

\noindent (a) 用主桡骨作响应变量对其他 4 种骨骼作回归.

\noindent (i) 提出并拟合适当的回归模型.

\noindent (ii) 分析残差.

\noindent (b) 用两种桡骨作响应变量对其他 3 种骨骼作多元多重回归.

\noindent (c) 对你在 (b) 中选择的模型和全模型计算 AIC.

\medskip
\noindent 7.23\quad 用\hyperref[tab:1.10]{表 1.10} 中在拍卖中出售的公牛特征数据:

\noindent (a) 作回归分析, 用 $Y_1=\mathrm{SalePr}$ 作响应变量, 预测变量为 Breed, YrHgt, FtFrBody, PrctFFB, Frame, BkFat, SaleHt 和 SaleWt.

\noindent (i) 通过只保留有单独显著性的预测变量来确定“最佳”回归方程.

\noindent (ii) 用最佳拟合模型, 对预测变量的一组值(按上面所列顺序) 5, 48.7, 990, 74.0, 7, 0.18, 54.2 和 1\,450, 构造销售价格的 95\% 预测区间.

\noindent (iii) 考察最佳拟合模型的残差.

\noindent (b) 用 $Y_1=\ln(\mathrm{SalePr})$ 作响应变量, 重复 (a) 中的分析. 你倾向于哪一种分析方法? 为什么?

\medskip
\noindent 7.24\quad 用\hyperref[tab:1.10]{表 1.10} 中拍卖出售公牛的特征数据:

\noindent (a) 进行回归分析, 响应变量为 $Y_1=\mathrm{SaleHt}$, 预测变量为 $Z_1=\mathrm{YrHgt}$ 和 $Z_2=\mathrm{FtFrBody}$.

\noindent (i) 拟合适当的模型并分析残差.

\noindent (ii) 当 $z_1=50.5$ 和 $z_2=970$ 时, 求 $Y_1=\mathrm{SaleHt}$ 的 95\% 预测区间.

\noindent (b) 以 $Y_1=\mathrm{SaleHt}$ 和 $Y_2=\mathrm{SaleWt}$ 作响应变量, $Z_1=\mathrm{YrHgt}$ 和 $Z_2=\mathrm{FtFrBody}$ 作预测变量, 进行多元回归分析.

\noindent (i) 拟合适当的多元模型并分析残差.

\noindent (ii) 对 $z_1=50.5$ 和 $z_2=970$, 求 $Y_1$ 和 $Y_2$ 的 95\% 预测椭圆. 将此椭圆与 (a)(ii) 中的预测区间作比较, 并加以评论.

\medskip
\noindent 7.25\quad 化学物质阿米替林(amitriptyline)被某些心理医生用作抗忧郁剂. 然而, 这种药物有一些副作用, 会引起心律不齐、血压异常和不规则的心电图波形等等. 对 17 个使用过量阿米替林的病人搜集了一组数据, 列于\hyperref[tab:7.6]{表 7.6}. 两个响应变量为

\noindent $Y_1=$ 总的 TCAD 血浆水平(TOT)

\noindent $Y_2=$ 在 TCAD 血浆水平中的阿米替林含量(AMI)

\noindent 五个预测变量为

\noindent $Z_1=$ 性别(GEN): 1 为女性, 0 为男性

\noindent $Z_2=$ 超剂量时抗忧郁剂的量(AMT)

\noindent $Z_3=$ PR 波测量值(PR)
